\documentclass[10pt,a4paper]{amsart}
\usepackage[margin=19mm]{geometry}
\usepackage{amsmath,amssymb,mathtools}
\usepackage[scr=rsfs,cal=euler]{mathalfa}
\usepackage{xfrac}
\usepackage[unicode,hidelinks,bookmarks=false]{hyperref}
\newcommand{\ie}{i.e.\ }
\newcommand{\cf}{cf.\ }
\newcommand{\resp}{respectively\ }
\newcommand{\Subsection}{\subsection}
\providecommand{\nobreakdash}{\nobreak\hbox{-}}
\setlength{\emergencystretch}{2em}
\multlinegap0pt
\usepackage{bm}
\usepackage{bbm}
\usepackage{dsfont}
\usepackage{xspace}
\usepackage{cancel}
\usepackage{mathtools}
\usepackage{cleveref}
\usepackage{amsthm}
\makeatletter
\@ifundefined{lemma}{\newtheorem{lemma}{Lemma}}{}
\@ifundefined{smallmatrix}{\newtheorem{smallmatrix}{Smallmatrix}}{}
\makeatother
\title[The eigenvalue decomposition of normal matrices by the skew-]{The eigenvalue decomposition of normal matrices by the skew-symmetric part}
\author{Simon Mataigne, Kyle A. Gallivan}
\date{Source: arXiv:2410.12421v4 (CC BY 4.0)}
\begin{document}
\maketitle

\noindent\emph{Source and licence.} The eigenvalue decomposition of normal matrices by the skew-symmetric part. Version arXiv:2410.12421v4 (CC BY 4.0), distributed under the Creative Commons Attribution 4.0 International licence, as recorded in the arXiv licence metadata for this exact version and shown on the abstract page https://arxiv.org/abs/2410.12421v4. Full rights notice: licenses/arxiv-2410.12421-CC-BY-4.0.txt.\par\smallskip

\noindent\textit{Editorial scope.} Counted unit: the Lemma whose source label is \texttt{lem:Invariance\_lambda} in the article (source0.tex lines 1456--1472, source label \texttt{lem:Invariance\_lambda}), delivered together with the article's own complete original proof of it (source0.tex lines 1473--1529, one \texttt{proof} environment of 57 lines). No mathematics is written, repaired or re-derived: statement and proof are the article's own text line for line. Required context retained uncounted: lem:permute (lines 1449--1452). Every definition, display and label of those lines is retained unchanged. Omitted from this package and not redistributed: 1--1448, 1453--1455, 1530--2478. Nothing inside a retained excerpt is omitted or altered: every retained body is delivered exactly as the article's lines are written, so no omission receipt of any kind accompanies this package. Citations: the retained excerpts carry no citation command at all, so no bibliography entry of the article is transcribed and no reference is redistributed. Reference adaptation: none was needed and none was made; every reference inside the delivered unit resolves inside it, so no unresolved reference or citation is produced. Printed slips kept literal: no printed word, symbol, hypothesis or number of the counted statement or of its proof was altered. Layout and class: the article's own class is \texttt{siamart1116} with packages including babel, graphicx, epstopdf, epsfig, amsfonts, fancyhdr, graphics, hyperref, amsmath, amssymb, colortbl, mathtools, cleveref, comment; the wrapper is the lane's pinned \texttt{amsart} editorial wrapper at 10pt on A4, which restates the theorem-like environments the retained text needs and adds the article's own macro definitions that the retained text needs (none), with extra wrapper packages (bm, bbm, dsfont, xspace, cancel, mathtools, cleveref). The delivered numbering is the unit's own. Assets: no retained span contains a figure environment, a graphics-inclusion command, a PDF-page inclusion, a table or a TikZ picture; no image file is copied into this package, the package declares no asset dependency and no image file is redistributed with it. Licence: version arXiv:2410.12421v4 of this article is distributed under the Creative Commons Attribution 4.0 International licence, as recorded in the arXiv licence metadata for this exact version and shown on the abstract page https://arxiv.org/abs/2410.12421v4. A full rights notice is delivered at licenses/arxiv-2410.12421-CC-BY-4.0.txt. Characters: every retained excerpt is pure ASCII, so the delivered engine, XeLaTeX with its default Latin Modern text font, reports no missing character.
\begin{lemma}\label{lem:permute}

	For every $A\in\mathbb{R}^{n\times n}$, all $Q\in\mathrm{ig}(A)$ and all $V\in\mathrm{O}(n)$, we have $VQV^\top \in\mathrm{ig}(VAV^\top )$. In particular, $V$ can be a permutation matrix.
	\end{lemma}

	\begin{lemma}
	\label{lem:Invariance_lambda}
		Let $\Sigma\in\mathbb{R}^{p\times p}$ be a diagonal matrix with distinct non-zero diagonal values. An orthogonal matrix $Q\in\mathrm{O}(2p)$ is such that
		\begin{equation}\label{eq:invariance_sigma}
			Q\begin{bmatrix}
			0&-\Sigma\\
			\Sigma& 0
			\end{bmatrix}Q^\top  = \begin{bmatrix}
			0&-\Sigma\\
			\Sigma& 0
			\end{bmatrix},
		\end{equation}
		if and only if $Q = \begin{bmatrix}
		\cos(\Phi)&-\sin(\Phi)\\
		\sin(\Phi)&\cos(\Phi)
		\end{bmatrix}$ for some diagonal matrix $\Phi\in\mathbb{R}^{p\times p}$.
	\end{lemma}

	\begin{proof}
	We proceed by induction. Assume first $p=1$, i.e., $\Sigma=\sigma \in \mathbb{R}$. Then, for every $Q\in\mathrm{O}(2)$, there is $\phi\in\mathbb{R}$ such that either  $Q=\left[\begin{smallmatrix}
		\cos(\phi)&-\sin(\phi)\\
		\sin(\phi)&\cos(\phi)
		\end{smallmatrix}\right]$ ($Q$ is a rotation) or $Q=\left[\begin{smallmatrix}
		\cos(\phi)&\sin(\phi)\\
		\sin(\phi)&-\cos(\phi)
		\end{smallmatrix}\right]$ ($Q$ is a reflection). It is easily verified that only rotations belong to $\mathrm{ig}\left(\left[\begin{smallmatrix}
		0&-\sigma\\
		\sigma&0
		\end{smallmatrix}\right]\right)$. Thus, \cref{lem:Invariance_lambda} holds for $p=1$. 
		
		Assume \cref{lem:Invariance_lambda} holds for $p\geq 1$ and let $\Sigma\in\mathbb{R}^{p\times p}$, $S_1 = \left[\begin{smallmatrix}
			0&-\Sigma\\
			\Sigma& 0
			\end{smallmatrix}\right]$ and $S_2 = \left[\begin{smallmatrix}
			0&-\sigma\\
			\sigma& 0
			\end{smallmatrix}\right]$ with $\sigma\in\mathbb{R}$ be non-zero and distinct from all entries of $\Sigma$. Let $Q=\left[\begin{smallmatrix}
			Q_1&Q_2\\
			Q_3&Q_4
			\end{smallmatrix}\right]\in\mathrm{O}(2p+2)$ with $Q_1\in\mathbb{R}^{2p\times 2p}$ and $Q_4\in\mathbb{R}^{2\times 2}$. Then, if one imposes\begin{equation*}
			Q\begin{bmatrix}
			S_1&0\\
			0&S_2
			\end{bmatrix}Q^\top =\begin{bmatrix}
			S_1&0\\
			0&S_2
			\end{bmatrix} \text{ and } Q\in \mathrm{O}(2p+2),
			\end{equation*}
			it must hold in particular that
			\begin{equation*}
			\begin{cases}
			Q_3 S_1 Q_3^\top  + Q_4 S_2 Q_4^\top  = S_2,\\
			Q_3 S_1 Q_1^\top  + Q_4 S_2 Q_2^\top  = 0,\\
			\end{cases} \text{ and } \begin{cases}
			Q_1^\top Q_1 + Q_3^\top Q_3 = I_{2p},\\
			Q_1^\top Q_2 + Q_3^\top Q_4 = 0.
			\end{cases}
			\end{equation*}
			By multiplying the second equation by $Q_1$ on the right, we obtain $Q_3 S_1 Q_1^\top Q_1 + Q_4 S_2 Q_2^\top Q_1 = 0$. Leveraging both orthogonality constraints, we obtain $Q_3S_1(I_{2p}-Q_3^\top Q_3)-Q_4S_2Q_4^\top Q_3=0$. Introducing the first equation, we obtain $Q_3S_1=S_2Q_3$. Let $Q_3:=[Q_{31}\ Q_{32}]$ such that \begin{align*}
		Q_3S_1=S_2Q_3 \iff Q_{32}\Sigma = S_2 Q_{31} &\text{ and } \ -Q_{31}\Sigma= S_2Q_{32}.\\
		\iff Q_{31}\Sigma^2 = \sigma^2 Q_{31} &\text{ and } \ Q_{32}\Sigma^2 = \sigma^2 Q_{32}.
\end{align*}			 
 Since $\sigma$ is distinct from all entries of $\Sigma$, $Q_3=0$. Hence, $Q_4\in\mathrm{O}(2)$ and $Q_2=0$. Finally, $Q\in \mathrm{ig}\left(\left[\begin{smallmatrix}
			S_1&0\\
			0&S_2
			\end{smallmatrix}\right]\right)$ if and only if it is block-diagonal, $Q_1\in\mathrm{ig}(S_1)$ and $Q_4\in\mathrm{ig}(S_2)=\mathrm{SO}(2)$.
			
			To conclude, we can define $\widetilde{\Sigma}:=\mathrm{diag}(\Sigma,\sigma)$ and take a permutation matrix $P$ such that $P\left[\begin{smallmatrix}
			S_1&0\\
			0&S_2
			\end{smallmatrix}\right]P^\top =\left[\begin{smallmatrix}
			0&-\widetilde{\Sigma}\\
			\widetilde{\Sigma}&0
			\end{smallmatrix}\right]$. Applying \cref{lem:permute} is enough to conclude that \cref{lem:Invariance_lambda} holds for $p+1$. Therefore, it holds for all $p\geq 1$.
	\end{proof}

\end{document}
